% ch06 —— §6 Adding masses（原文 wcf-paper.tex 3002–3152 行）
\section{加入质量}
\label{sec:masses}

本节我们简要说明：若要纳入非平凡的质量参数，前几节的结果应当如何修改。

\subsection{含质量的单粒子修正}

第 4 节所考虑的 $U(1)$ 理论有一个简单的变体：我们可以考虑含有若干超多重态的 $U(1)$ 理论，这些超多重态带电，电荷为 $q_i$。% TODO 合稿改 \ref{sec:periodicnut}
粒子种类多于一种的理论会牵涉味荷，并且会非平凡地依赖于质量参数。四维中的质量参数是复数 $m_i$。紧致化到三维后，会多出一个实的周期质量参数 $m^3_i$，它本质上是味对称的一条 Wilson 线。我们记 $\psi_i := 2\pi R m^3_i$，其周期为 $2\pi$。

质量参数以非常简单的方式进入修正后的度量：每个粒子对 $V$ 与 $A$ 给出一个可加贡献，其形式与我们此前遇到的那个相仿，
\begin{equation} \label{eq:V-sum-masses}
V = \sum_i \frac{q_i^2 R}{4 \pi} \sum_{n=-\infty}^\infty \left(
\frac{1}{\sqrt{R^2 \abs{q_i a + m}^2 + (q \frac{\theta_\elec}{2 \pi} +
\frac{\psi_i}{2 \pi} + n)^2}} - \kappa_n \right)
\end{equation}

坐标 $\CX_\elec$ 不变：
\begin{equation}
\CX_\elec = \exp \left[ \pi R \frac{a}{\zeta} + i \theta_\elec +
\pi R \zeta \bar{a} \right].
\end{equation}
引入质量参数的类似组合也很有用：
\begin{equation}
\mu_i := \exp \left[ \pi R \frac{m_i}{\zeta} + i \psi_i + \pi R
\zeta \bar{m_i} \right].
\end{equation}
半平坦的 $\CX_\magn$ 含有把 $d=4$ 中所有粒子积分掉所产生的贡献：
\begin{equation}
\begin{aligned}
\CX_\magn^{\sfm}(\zeta) &= e^{i \theta_\magn} \times \\
&\quad \prod_i \exp\left[-i\frac{R q_i}{2
\zeta} \left((q_i a+m_i) \log\frac{q_i a+ m_i}{e \Lambda} \right) +
i \frac{\zeta R q_i}{2} \left((q_i \bar a + \bar
m_i)\log\frac{q_i \bar a+ \bar m_i}{e\bar \Lambda} \right)\right].
\end{aligned}
\end{equation}
$\CX_\magn^{\sfm}(\zeta)$ 绕 $q_i a+m_i =0$ 的单值变换为
\begin{equation}\label{eq:mono2-masses}
\CX_\magn^{\sfm} \to (-\mu_i)^{q_i} \CX_\elec^{q_i^2}
\CX_\magn^{\sfm}.
\end{equation}
完整的坐标同样含有来自所有这些粒子的瞬子贡献，
\begin{equation}\label{eq:FullX2-masses}
\begin{aligned}
\CX_\magn = \CX_\magn^{\sfm} \prod_i \exp\Biggl[ & \frac{i q_i}{4\pi}
\int_{\ell^i_+} \frac{d \zeta'}{\zeta'}
\frac{\zeta' + \zeta}{\zeta' - \zeta} \log[1-\mu_i \CX_\elec(\zeta')^{q_i}] \\
& - \frac{i q_i}{4\pi}\int_{\ell^i_-}
\frac{d \zeta'}{\zeta'} \frac{\zeta' + \zeta}{\zeta' - \zeta}
\log[1-\mu_i^{-1} \CX_\elec(\zeta')^{-q_i}] \Biggr],
\end{aligned}
\end{equation}
其中围道 $\ell^i_\pm$ 取为 $\zeta$ 平面中连接 $0$ 与 $\infty$、且落在下面两个半平面内的任意路径：
\begin{equation}
\CU^i_\pm = \left\{ \zeta: \pm {\rm Re} \frac{q_i a + m_i}{\zeta} < 0 \right\}.
\end{equation}

不连续性如今依赖于质量：
\begin{subequations} \label{eq:X2-disc-masses}
\begin{align}
(\CX_\magn)_{\ell^i_+}^{+} &= (\CX_\magn)_{\ell^i_+}^{-} (1 - \mu_i \CX_\elec^q)^{-q}, \\
(\CX_\magn)_{\ell^i_-}^{+} &= (\CX_\magn)_{\ell^i_-}^{-} (1 - \mu_i^{-1}
\CX_\elec^{-q})^q.
\end{align}
\end{subequations}

本节的所有公式也都可以沿着第 4.7 节的思路推广到更高秩。% TODO 合稿改 \ref{sec:higher-rank}

\subsection{含质量的多粒子修正}

现在我们已有条件理解质量参数在一般的黎曼--希尔伯特问题与微分问题中所起的作用。考虑一个具有 $n_f$ 个味对称性的规范理论。把味荷记为 $\gamma^f$，并由质量与味 Wilson 线以显然的方式构造出 $\mu_{\gamma^f}$：
\begin{equation}
\mu_{\gamma^f} := \exp \left[ \pi R \frac{m_{\gamma^f}}{\zeta} + i \psi_{\gamma^f} + \pi R
\zeta \bar{m}_{\gamma^f} \right].
\end{equation}
阿贝尔问题的不连续性提示出如下形式的广义 Kontsevich--Soibelman 因子：
\begin{equation}
\label{eq:KS-generalized-masses}
\CK_{\gamma,\gamma^f} := \CX_{\gamma'} \to \CX_{\gamma'} (1-
\sigma(\gamma) \mu_{\gamma^f} \CX_\gamma)^{\langle \gamma,\gamma'
\rangle}.
\end{equation}
于是我们可以定义一个与第 5.1 节相仿的黎曼--希尔伯特问题，它对每个电荷为 $\gamma$、味荷为 $\gamma^f$ 的粒子配上不连续性 $\CK_{\gamma,\gamma^f}$。% TODO 合稿改 \ref{sec:defining-rh}
假定推广到使用这些修改后的辛同胚之后，穿墙公式仍然给出正确的 BPS 简并度，我们就可以用这个黎曼--希尔伯特问题在 $\CM$ 上构造一个超凯勒度量；我们提出，它就是正确的那个度量。

超对称场论中有一个标准技巧：把质量参数视为某个扩充理论的矢量多重态标量的真空期望值，而在该扩充理论中，味对称性已被弱规范化 \cite{Seiberg:1993vc}。由于是弱规范化，带磁味荷的粒子非常重，在味规范耦合趋于零的极限下可以略去；此时剩下的味规范荷只有电性的。利用这个技巧，含质量的广义穿墙公式可以解释为标准穿墙公式在零耦合极限下的情形。

最后，我们希望把第 5.5 节的微分表述推广到能处理质量参数。% TODO 合稿改 \ref{sec:differential-problem}
在那一节中，我们依赖于诸因子 $S_\ell$ 与参数无关这一事实。在修改后的问题中，诸 $S_\ell$ 显式依赖于 $\mu_{\gamma^f}$，从而依赖于 $m_i$、$R$、$\zeta$。然而，所有 $S_\ell$ 确实都被 $\partial_{m_i} + i \frac{\pi R}{\zeta} \partial_{\psi_i}$ 零化，也被 $\partial_{\bar m_i} + i \pi R \zeta \partial_{\psi_i}$ 零化。于是把那一节的论证稍作修改，即可说明：黎曼--希尔伯特问题的解 $\CX$ 满足如下形式的微分方程
\begin{align}
\partial_{m_i} \CX &= \left(\frac{1}{\zeta} \CA^{(-1)}_{m_i} + \CA^{(0)}_{m_i} \right) \CX,
\label{eq:mass-equation} \\
\partial_{\bar{m}_i} \CX &= \left(
\CA^{(0)}_{\bar{m}_i} + \zeta \CA^{(1)}_{\bar{m}_i} \right) \CX, \label{eq:mass-bar-equation}
\end{align}
其中 $\CA^{(-1)}_{m_i}$ 不仅是 $\CM_u$ 上的矢量场，还包含算子 $-i \frac{\pi R}{\zeta} \partial_{\psi_i}$；对 $\CA^{(1)}_{\bar{m}_i}$ 也有类似的说法。

诸 $S_\ell$ 也被 R 对称性与标度不变性的算子零化，因此我们得到 \eqref{eq:rescaling-m1}、\eqref{eq:rescaling-m2} 的类似物（在转到适当的规范之后），
\begin{align}
\left(R \partial_R - a^I \partial_{a^I} - \bar{a}^I \partial_{\bar{a}^I} - \Lambda \partial_\Lambda - \bar{\Lambda} \partial_{\bar{\Lambda}} - m^i \partial_{m^i} - \bar{m}^i \partial_{\bar{m}^i} \right) \CX &= 0, \label{eq:rescaling-m1} \\
\left(\zeta \partial_\zeta + a^I \partial_{a^I} - \bar{a}^I \partial_{\bar{a}^I} + \Lambda \partial_\Lambda - \bar{\Lambda} \partial_{\bar{\Lambda}} + m^i \partial_{m^i} - \bar{m}^i \partial_{\bar{m}^i} \right) \CX &= 0. \label{eq:rescaling-m2}
\end{align}
利用这些方程，我们可以得到 $\CX$ 对 $R$ 与 $\zeta$ 的依赖所满足的标准形式 (5.34)、(5.35)，同样带有这样的修改：$\CA_R$ 与 $\CA_\zeta$ 现在包含对诸 $\psi_i$ 的导数。% TODO 合稿改 \eqref{eq:R-equation}、\eqref{eq:zeta-equation}

\begin{yizhu}
\begin{enumerate}[leftmargin=2em,itemsep=2pt,topsep=2pt]
\item 式 \eqref{eq:V-sum-masses} 中的 $\abs{q_i a + m}$ 与 $q\,\theta_\elec/(2\pi)$，原文分别印作 $m$ 与 $q$（不带下标 $i$），而下文同类式子一律带下标；译文公式照原文逐字保留，此处仅提示读者：这两处应按 $m_i$、$q_i$ 理解。
\item 式 \eqref{eq:KS-generalized-masses} 中的 $\sigma(\gamma)\in\IZ_2$ 是第 2.2 节引入的二次加细（quadratic refinement），满足 $\sigma(\gamma_1)\sigma(\gamma_2)=(-1)^{\langle\gamma_1,\gamma_2\rangle}\sigma(\gamma_1+\gamma_2)$；本块沿用它，不另作定义。% TODO 合稿改 \ref{sec:ReviewKS}
\end{enumerate}
\end{yizhu}
